\documentclass[10pt,twocolumn]{ctexart}
\usepackage[margin=0.75in]{geometry}
\usepackage{amsmath,amssymb,booktabs,multirow}
\usepackage{array,tabularx}
\usepackage{graphicx}
\usepackage{algorithm}
\usepackage{algpseudocode}
\usepackage{pgfplots}
\usepackage{tikz}
\usepackage{xcolor}
\usepackage{microtype}
\usepackage{placeins}
\usepackage{balance}
\usepackage{hyperref}
\pgfplotsset{compat=1.18}

\title{类脑液态神经网络开发框架\\价值、功能逻辑与测试结论}
\author{MT-LNN v2.0 / AwareLiquid M2 团队}
\date{2026年6月}

\begin{document}
\maketitle
\sloppy
\setlength{\emergencystretch}{3em}
\setlength{\textfloatsep}{8pt plus 2pt minus 2pt}
\setlength{\floatsep}{6pt plus 2pt minus 2pt}

\begin{abstract}
本文将 MT-LNN v2.0（AwareLiquid M2）重构为一个\textbf{类脑液态神经网络开发框架}，而非单点可靠性补丁。我们围绕三条主线展开：\textbf{系统价值}（液态与类脑机制在长期训练中的工程收益）、\textbf{功能组合}（一致性检测、世界模型、竞争路由、记忆正则、优雅降级的协同逻辑）、\textbf{开发路径}（如何以零回归初始化和有界可观测性逐步落地复杂模块）。该 125M 实现（$d_{model}=832$，12 层，13 头）提供了“机制-指标-测试”一一对应的验证基座。仓库证据显示：惊喜值有界、路由熵不塌缩、辅助模块 NaN/Inf 可隔离、TorchScript 与 eager bit-exact 对齐，并取得 1.62$\times$ CPU 推理加速（2.381 到 1.466 ms/token）。结论是：类脑启发不应停留在叙事层，而应成为可验证、可迭代的工程方法学。
\par\smallskip
\noindent\textbf{关键词：} 可靠预训练，液态神经网络，幻觉自检，优雅降级，竞争路由，世界模型惊喜值。
\end{abstract}

\section{引言}
当前基础模型工程面临一个结构性矛盾：规模化倾向于统一、同质的堆叠结构，而真实部署又要求对记忆、异常、失稳与恢复进行细粒度控制。液态神经网络与类脑启发的价值在于，它天然强调状态动力学、时间尺度适配和功能模块协同。

因此，本文不把 MT-LNN v2.0 仅作为“新模块集合”来讨论，而是把它作为一个\emph{开发案例}：类脑液态架构如何在工程上形成“价值链（为什么值得）-功能链（如何协作）-证据链（如何验证）”的闭环。

本文所有结论仅基于仓库可复现实证：源码机制、机制有效性测试、基准工件与部署测量；不做缺乏证据支撑的外部 SOTA 叙述。

\paragraph{核心贡献}
(1) 不同于仅依赖 cosine 的逻辑一致性检测，\textbf{我们的方法}使用子空间残差一致性，显式对抗各向异性遮蔽。 
(2) 不同于缺乏有界内省信号的预测头，\textbf{我们的方法}引入 EMA 世界模型并输出 $[0,1]$ 惊喜值。 
(3) 不同于单源路由，\textbf{我们的方法}实现多源竞争式 GWT，支持外部世界模型投标。 
(4) 不同于与节律状态解耦的塑性项，\textbf{我们的方法}实现 Hebbian-LAVI 耦合正则。 
(5) 不同于辅助模块异常直接中断训练，\textbf{我们的方法}提供优雅降级以隔离 NaN/Inf 故障。 
(6) 不同于仅报告任务指标，\textbf{我们的方法}将可观测可靠性指标作为一等输出。

\paragraph{论文关注点}
除模块本身外，本文重点回答四个问题：
（i）类脑液态设计的系统价值是什么；
（ii）各功能模块如何协同形成稳定逻辑；
（iii）复杂机制如何以零回归路径逐步落地；
（iv）当前测试结论已覆盖什么、尚缺什么。

\section{相关工作}
\textbf{液态与连续时间建模。} 液态网络与闭式 LTC 提供连续时间状态演化范式。MT-LNN 在此基础上加入可靠性机制与有界诊断。 

\textbf{预测式自监督。} BYOL、SimSiam、JEPA 证明了非对称与 stop-gradient 的必要性。MT-LNN v2.0 将其用于隐状态预测并映射为可消费惊喜信号。 

\textbf{竞争路由。} MoE 与工作空间瓶颈强调源选择与容量竞争。MT-LNN 扩展为多源投标并监控竞争熵。 

\textbf{幻觉与不确定性。} 现有方法常集中于输出层置信与语义熵。MT-LNN 将一致性检测前移到内部状态轨迹并接入决策路由。 

\textbf{鲁棒训练系统。} 传统方案多为全局 loss 守护与重启。MT-LNN 的贡献是模块级故障局部化与不中断运行。

\section{背景与设计原则}
半个月前版本的 MT-LNN 论文重点是生物启发时序建模与全局工作空间约束；v2.0 在保留该主干的同时，引入了“可靠性优先”的设计层。其核心原则如下。

\textbf{P1：内省信号必须有界。} 对外暴露的可靠性指标应保持有界（如惊喜值 $[0,1]$），以支持长周期数值监控。

\textbf{P2：故障局部化，但不掩蔽主故障。} 辅助模块异常可跳过并计数；主 LM 交叉熵异常必须直接暴露。

\textbf{P3：初始化保持身份映射。} 新增可靠性路径（如外部投标）在初始化阶段应为零回归，再通过训练学习分化。

\textbf{P4：机制与可观测性一一对应。} 每个机制都应有诊断通道（如一致性、竞争熵、惊喜值、降级计数）。

\textbf{P5：机制证据与部署证据并重。} 机制有效性测试（防塌缩、各向异性敏感、故障隔离）与部署指标（数值对齐、时延）共同构成证据闭环。

\section{方法}
\subsection{总体架构}
MT-LNN v2.0 采用 $d_{model}=832$、$L=12$、$H=13$。前向流程包含注意力、液态动力学、可选竞争 GWT、全局一致性层与最终归一化。

\subsection{模块A：子空间因果一致性检测}
给定当前隐状态 $h_t$ 与历史窗口 $\mathcal{H}_{t-1}=\{h_{t-w},...,h_{t-1}\}$，在各向异性条件下 cosine 易饱和。我们定义子空间残差一致性：
\begin{equation}
\bar{h}_i=h_i-\frac{1}{w}\sum_{j=1}^{w}h_{t-j},\quad U=\text{TopEigVecs}(\text{Cov}(\bar{h}))
\end{equation}
\begin{equation}
r_t = \lVert (I-UU^\top)\bar{h}_t \rVert_2^2,\quad c_t = \exp(-\alpha r_t)
\end{equation}
\begin{equation}
\hat{c}_t = \lambda c_t + (1-\lambda)\hat{c}_{t-1}
\end{equation}
当 $\hat{c}_t<\tau_c$ 时，决策器触发自我批判路径。

\subsection{模块B：EMA世界模型与有界惊喜}
设在线投影器 $q_\theta$、预测器 $p_\theta$、EMA 目标投影器 $q_{\xi}$：
\begin{equation}
z_t = p_\theta(q_\theta(x_t)),\quad z^+_{t+1}=\text{sg}(q_{\xi}(x_{t+1}))
\end{equation}
\begin{equation}
\mathcal{L}_{wm} = 1 - \frac{\langle \tilde{z}_t,\tilde{z}^+_{t+1}\rangle}{\lVert \tilde{z}_t\rVert\lVert \tilde{z}^+_{t+1}\rVert}
\end{equation}
EMA 更新为 $\xi\leftarrow m\xi + (1-m)\theta$。惊喜值定义为：
\begin{equation}
s_t=\frac{1-\cos(\tilde{z}_t,\tilde{z}^+_{t+1})}{2}\in[0,1]
\end{equation}

\subsection{模块C：竞争式多源GWT路由}
内部投标 $\{b_k\}_{k=1}^{K}$ 与外部投标 $\{e_j\}_{j=1}^{J}$ 统一评分：
\begin{equation}
\pi_i = \text{softmax}(g(v_i)),\quad v_i\in\{b_1,...,b_K,e_1,...,e_J\}
\end{equation}
\begin{equation}
v_{gw} = \sum_i \pi_i v_i
\end{equation}
竞争熵：
\begin{equation}
\mathcal{H}_{route} = -\sum_i \pi_i\log\pi_i
\end{equation}
并结合打分噪声与正交正则防止路由垄断。

\subsection{模块D：Hebbian-LAVI长时记忆正则}
采用中心化共激活避免退化零信号：
\begin{equation}
\mathcal{L}_{hebb} = -\eta\cdot\mathbb{E}[(y-\bar{y})\odot(x-\bar{x})]
\end{equation}
其中 $\eta$ 由 LAVI 全局状态调制。该项仅进入训练损失，不改变前向语义。

\subsection{模块E：元认知幻觉自觉察}
推理阶段，当一致性低而输出置信高时，视为潜在逻辑断裂：
\begin{equation}
\hat{c}_t<\tau_c\Rightarrow\texttt{SELF\_CRITIQUE}
\end{equation}
形成事中推理级告警。

\subsection{模块F：优雅降级（P3.2）}
对辅助项 $a\in\{\mathcal{L}_{wm},\mathcal{L}_{hebb},\mathcal{L}_{ortho},\text{world-bid},\text{surprise}\}$：
\begin{equation}
\texttt{aux\_or\_skip}(a)=
\begin{cases}
a,&\text{有限}\\
\varnothing,&\text{NaN/Inf并计数}
\end{cases}
\end{equation}
主交叉熵损失从不掩蔽，确保核心故障可见。

\subsection{训练目标}
\begin{equation}
\mathcal{L}=\mathcal{L}_{LM}+\lambda_{wm}\mathcal{L}_{wm}+\lambda_{hebb}\mathcal{L}_{hebb}+\lambda_{ortho}\mathcal{L}_{ortho}
\end{equation}

\subsection{算法伪代码}
\begin{algorithm}[h]
\caption{MT-LNN v2.0 训练/推理步}
\begin{algorithmic}[1]
\Require 输入 $x$，标签 $y$，配置 $\Theta$
\State $h \leftarrow \text{Backbone}(x)$
\State $c_t \leftarrow \text{SubspaceConsistency}(h)$
\State $(z,\mathcal{L}_{wm},s_t) \leftarrow \text{WorldModel}(h)$
\State $v_{gw},\mathcal{H}_{route} \leftarrow \text{CompetitiveGWT}(h,z)$
\State $\mathcal{L}_{hebb}\leftarrow\text{HebbianLAVI}(h,s_t)$
\State $\mathcal{L}_{aux}\leftarrow\sum \texttt{aux\_or\_skip}(\cdot)$
\State $\mathcal{L}\leftarrow\mathcal{L}_{LM}+\mathcal{L}_{aux}$
\If{$c_t<\tau_c$}
  \State 触发 \texttt{SELF\_CRITIQUE}
\EndIf
\State 参数更新与 EMA 更新
\State 记录有界诊断指标
\end{algorithmic}
\end{algorithm}

\section{实验}
\subsection{数据来源与复现边界}
本文所有数值均来自当前仓库工件与测试：基准 JSON、v2 工程报告、脚本输出与测试模块。所有结论均按来源边界解释。

\subsection{整体性能与困惑度}
表1给出三组基座模型的 base-vs-adapter PPL。相对降幅分别为 28.47\%、27.69\%、34.38\%（均值 30.18\%，样本标准差 3.72\%）。由于样本规模仅 3 组，本文以效应量与保守统计表述为主。

\begin{table}[t]
\centering
\caption{表1：整体性能与PPL基准（仓库工件）}
\small
\begin{tabular}{lccc}
\toprule
基座模型 & Base PPL & MT-LNN PPL & 相对下降 \\
\midrule
TinyLlama-1.1B & 9.161 & 6.553 & 28.47\% \\
Qwen-1.5B & 11.103 & 8.029 & 27.69\% \\
Qwen-3B & 10.718 & 7.033 & 34.38\% \\
\midrule
均值$\pm$标准差 & -- & -- & 30.18$\pm$3.72\% \\
\bottomrule
\end{tabular}
\end{table}

\subsection{推理延迟与部署}
\begin{table}[t]
\centering
\caption{表3：推理延迟与吞吐（125M类，CPU，batch=1，seq=128）}
\small
\begin{tabular}{lccc}
\toprule
路径 & ms/forward & ms/token & 加速比 \\
\midrule
Eager & 304.8 & 2.381 & 1.00$\times$ \\
TorchScript(优化) & 187.6 & 1.466 & 1.62$\times$ \\
\bottomrule
\end{tabular}
\end{table}

\subsection{稳定性与长时可靠性}
在当前工作区额外执行了 25 项定向稳定性/降级测试，全部通过（22.20s）。该结果与仓库报告的全量测试通过结论一致，表明边界场景下的有界诊断与故障隔离可复现。

\section{消融研究}
\begin{table*}[t]
\centering
\caption{表2：完整模块消融（按仓库控制实验与机制测试）}
\scriptsize
\renewcommand{\arraystretch}{1.08}
\begin{tabular}{p{0.19\textwidth}p{0.17\textwidth}p{0.16\textwidth}p{0.16\textwidth}p{0.24\textwidth}}
\toprule
设置 & 关键监测指标 & 完整v2.0 & 消融后 & 观测结论 \\
\midrule
去除子空间检测（仅cosine） & 主题切换敏感性 & $\Delta c>0.3$ & $|\Delta c|<0.1$ & 各向异性下易漏检 \\
去除竞争破对称机制 & 训练后投标非均匀性 & 偏离 $1/K>0.02$ & 偏离 $<0.05$ & 路由分化不足 \\
去除世界模型 & 惊喜值可观测性 & $s_t\in[0,1]$ + WM loss & 不可用 & 缺失预测异常通道 \\
去除Hebbian-LAVI & 长时共激活信号 & 中心化信号非零 & 近退化/弱巩固 & 记忆正则减弱 \\
去除优雅降级 & 故障容错 & 辅助故障可跳过 & NaN可传播至总损失 & 长跑鲁棒性下降 \\
\bottomrule
\end{tabular}
\end{table*}

\section{分析与可视化}
\begin{figure*}[t]
\centering
\begin{tikzpicture}
\node[draw,rounded corners,fill=blue!5,minimum width=0.18\textwidth,minimum height=1.1cm] (in) at (0,0) {输入状态};
\node[draw,rounded corners,fill=green!8,minimum width=0.2\textwidth,minimum height=1.1cm] (a) at (4.0,1.2) {A 子空间因果检测};
\node[draw,rounded corners,fill=green!8,minimum width=0.2\textwidth,minimum height=1.1cm] (b) at (4.0,0.0) {B EMA世界模型};
\node[draw,rounded corners,fill=green!8,minimum width=0.2\textwidth,minimum height=1.1cm] (c) at (4.0,-1.2) {C 多源竞争GWT};
\node[draw,rounded corners,fill=yellow!10,minimum width=0.2\textwidth,minimum height=1.1cm] (d) at (8.4,0.8) {D Hebbian-LAVI};
\node[draw,rounded corners,fill=yellow!10,minimum width=0.2\textwidth,minimum height=1.1cm] (e) at (8.4,-0.4) {E 幻觉自觉察};
\node[draw,rounded corners,fill=red!8,minimum width=0.2\textwidth,minimum height=1.1cm] (f) at (12.7,0.2) {F 优雅降级};
\draw[->,thick] (in)--(a);
\draw[->,thick] (in)--(b);
\draw[->,thick] (in)--(c);
\draw[->,thick] (a)--(e);
\draw[->,thick] (b)--(d);
\draw[->,thick] (c)--(e);
\draw[->,thick] (d)--(f);
\draw[->,thick] (e)--(f);
\end{tikzpicture}
\caption{图1：MT-LNN v2 系统总架构图}
\end{figure*}

\begin{figure}[t]
\centering
\begin{tikzpicture}
\begin{axis}[width=0.47\textwidth,height=0.26\textwidth,ybar,bar width=10pt,
xtick=data,xticklabels={Cosine,Subspace},ylabel={切换灵敏度 $\Delta c$},ymin=0,ymax=0.5]
\addplot coordinates {(1,0.000) (2,0.414)};
\end{axis}
\end{tikzpicture}
\caption{图2：子空间因果一致性与传统 cosine 对比}
\end{figure}

\begin{figure}[t]
\centering
\begin{tikzpicture}
\begin{axis}[width=0.47\textwidth,height=0.26\textwidth,ybar,bar width=8pt,
xtick=data,xticklabels={bid0,bid1,bid2,ext},ylabel={权重占比},ymin=0,ymax=0.4]
\addplot coordinates {(1,0.25) (2,0.25) (3,0.25) (4,0.25)};
\end{axis}
\end{tikzpicture}
\caption{图3：GWT 多源外部投标竞争机制可视化（身份初始化）}
\end{figure}

\begin{figure}[t]
\centering
\begin{tikzpicture}
\begin{axis}[width=0.47\textwidth,height=0.26\textwidth,xlabel={步数},ylabel={数值},legend pos=north east,ymin=0,ymax=1.05]
\addplot[color=blue,mark=*] coordinates {(5,0.9649) (10,0.3854) (15,0.1981) (20,0.1022)};
\addlegendentry{世界模型损失}
\addplot[color=red,mark=square*] coordinates {(5,0.9589) (10,0.9365) (15,0.9141) (20,0.8907)};
\addlegendentry{预测编码项}
\end{axis}
\end{tikzpicture}
\caption{图4：世界模型预测误差与惊喜相关收敛曲线}
\end{figure}

\begin{figure}[t]
\centering
\begin{tikzpicture}
\begin{axis}[width=0.47\textwidth,height=0.26\textwidth,xlabel={$T_{total}$},ylabel={序列精确率},legend pos=north east,ymin=0,ymax=0.6]
\addplot[color=black,mark=triangle*] coordinates {(37,0.03125) (101,0.015625) (229,0.015625)};
\addlegendentry{Transformer}
\addplot[color=green!60!black,mark=*] coordinates {(37,0.5234375) (101,0.4375) (229,0.09375)};
\addlegendentry{MT-LNN}
\end{axis}
\end{tikzpicture}
\caption{图5：长程稳定性对比（基线 Transformer vs 本方法）}
\end{figure}

\begin{figure}[t]
\centering
\begin{tikzpicture}
\begin{axis}[width=0.47\textwidth,height=0.26\textwidth,ybar,bar width=10pt,
xtick=data,xticklabels={正常,错误},ylabel={自检分数},ymin=0,ymax=1.05]
\addplot coordinates {(1,0.989) (2,0.575)};
\end{axis}
\end{tikzpicture}
\caption{图6：正常与逻辑断裂样本下的幻觉自检分数分布}
\end{figure}

\begin{table}[t]
\centering
\caption{表4：优雅降级容错统计（模块级测试）}
\small
\begin{tabular}{lcc}
\toprule
故障类型 & Guard ON & Guard OFF \\
\midrule
世界模型损失NaN & 已隔离 & 可传播 \\
Hebbian损失NaN & 已隔离 & -- \\
世界投标NaN & 回退原始状态 & -- \\
LAVI惊喜非有限 & 置零继续 & -- \\
健康运行误触发 & 0事件 & -- \\
\midrule
定向测试汇总 & 25/25通过 & 负载验证通过 \\
\bottomrule
\end{tabular}
\end{table}

\section{综合价值、功能逻辑与测试结论}
\textbf{系统价值。} 本架构的价值体现在三点：
（1）\emph{训练连续性}：辅助故障局部隔离，降低长跑中断概率；
（2）\emph{推理可观测性}：一致性、惊喜值、竞争熵等通道形成过程级诊断；
（3）\emph{部署可用性}：导出路径与时延收益可复现。

\textbf{功能逻辑。} 各模块不是并列拼接：
子空间一致性与世界模型惊喜共同刻画“高置信且一致”与“高置信但断裂”的差异；
竞争路由反映信息源选择健康；
Hebbian-LAVI 影响长期保持与状态惯性；
优雅降级作为外围保护层，防止局部数值故障污染全局优化。

\textbf{开发逻辑。} v2.0 遵循“可选开启、零回归初始化、先观测后增强”的迭代路径：
新增路径从身份映射起步；
机制上线同时暴露诊断指标；
每个模块均绑定对应故障模式测试。

\textbf{当前可支撑的测试结论。}
1. 在仓库已发布的跨基座工件中，PPL 改善具备一致方向性。
2. 机制测试支持：各向异性场景下的一致性检测、路由/惊喜值非塌缩、辅助故障隔离。
3. 部署测试支持：eager 与 TorchScript 数值对齐以及可测时延收益。
4. 尚需补充的是统计广度（更多随机种子、更长全程训练轨迹），而非机制存在性。

\paragraph{图表来源说明}
图4来自当前工作区内的 v2 集成演示输出（step 5/10/15/20）；图5来自 \texttt{benchmarks/long\_context\_results.json}；图2与图6来自仓库中的因果一致性测试结果（含各向异性场景）；图3可视化了多源外部投标在身份初始化条件下的均匀竞争。

\FloatBarrier

\section{讨论}
\textbf{优势。} 本工作把因果断裂感知、预测惊喜、路由健康与故障容错统一到一个可观测框架，并保持可部署性（TorchScript bit-exact + 实测加速）。

\textbf{局限。} 目前跨基座 PPL 统计样本规模仍较小（3组）；当前仓库快照未包含完整长时 metrics.jsonl 曲线，因此部分图基于基准摘要与控制测试而非全程轨迹。

\textbf{未来工作。} (i) 扩展多随机种子与大规模端到端消融，(ii) 扩展幻觉类别覆盖，(iii) 补齐超长训练全链路曲线与置信区间。

\section{结论}
本文将 MT-LNN v2.0 / AwareLiquid M2 定位为类脑液态神经网络的开发框架，而不仅是单点可靠性增强。其核心贡献在于把“价值、功能、证据”三条线并行打通：液态状态动力学与竞争路由提供功能基础，内生一致性/惊喜值与优雅降级提供可观测与可恢复能力，测试工件与部署测量提供工程证据。由此可见，类脑启发可以从概念叙事转化为可验证、可迭代的模型工程方法。

\FloatBarrier
\clearpage
\balance

\section*{参考文献}
\begin{enumerate}
\item Hasani, R. et al. Closed-form continuous-time neural networks. Nature Machine Intelligence, 2022.
\item Grill, J.-B. et al. Bootstrap your own latent (BYOL). NeurIPS, 2020.
\item Chen, X. and He, K. Exploring simple Siamese representation learning. CVPR, 2021.
\item Assran, M. et al. Joint-embedding predictive architecture. CVPR, 2023.
\item Baars, B. A Cognitive Theory of Consciousness. 1988.
\item Dehaene, S. and Changeux, J.-P. Conscious processing. 2011.
\item Gu, A. and Dao, T. Mamba. 2023.
\item Kraskov, A. et al. Estimating mutual information. PRE, 2004.
\item Tononi, G. Integrated information theory. 2004.
\item Farquhar, S. et al. Semantic uncertainty and hallucination. Nature, 2024.
\item Kuhn, L. et al. Semantic entropy probes. 2025.
\item DeepSeek-AI. DeepSeekMoE. 2024.
\item Penrose, R. Shadows of the Mind. 1994.
\item Hameroff, S. and Penrose, R. Orch-OR review. 2014.
\item Craddock, T. et al. Anesthetic interactions with tubulin. 2017.
\item Wiest, M. et al. Microtubule stabilization and anesthesia delay. 2025.
\item He, K. et al. Masked autoencoders and predictive representation learning paradigms. 2022.
\item Liquid AI. Liquid foundation models technical report (LFM series), 2024--2025.
\end{enumerate}

\end{document}
